\documentclass[../../../include/open-logic-section]{subfiles}

\begin{document}
	
\olfileid[tr]{sth}{z}{nat}
\olsection{Doğal Sayılarımızı Seçmek}

%Let us take it for granted that the informal principle \stagesinf{}
%can be justified. It is also, though, worth commenting on the
%particular axiom we extracted from that informal principle. 

\olref[infinity-again]{defnomega} içinde ``doğal sayılar'' ifadesini
açıkça tanımladık. Bu koşul koymayı nasıl anlamalısınız? Bu metafizik bir
iddia değil, yalnızca belirli kümeleri doğal sayılar \emph{olarak ele alma}
kararıdır. Böyle düşünmenin nedenlerine
\olref[sfr][rel][ref]{sec}, \olref[sfr][arith][ref]{sec} ve
\olref[sfr][infinite][dedekindsproof]{sec} içinde değindik. Ancak bu
nedenleri daha da belirginleştirebiliriz.

Sonsuzluk Aksiyomumuz \citet{VonNeumann1925} yaklaşımını izler. Bunun
yerine kabul edebileceğimiz başka bir aksiyom şudur:

\begin{defish}
\emph{Zermelo'nun \citeyear{Zermelo1908Untersuchungen} tarihli Sonsuzluk
Aksiyomu.} $\emptyset \in A$ olan ve
$(\forall x \in A)\{x\} \in A$ koşulunu sağlayan bir $A$ kümesi vardır.
\end{defish}

Bizim von Neumann esinli Sonsuzluk Aksiyomumuz yerine Zermelo'nun
aksiyomunu kullansaydık yine bir Dedekind-sonsuz küme, dolayısıyla bir
Dedekind cebiri elde ederdik. Zermelo'nun yaklaşımında cebirimizin ayırt
edilmiş elemanı yine $\emptyset$ ($0$ yerine kullandığımız nesne) olurdu;
fakat bire bir fonksiyon $x\mapsto x\cup\{x\}$ yerine
$x\mapsto\{x\}$ eşlemesiyle verilirdi. Bunun en basit sonucu, Zermelo'nun
$2$'yi $\{\{\emptyset\}\}$ olarak; bizimse (von Neumann ile birlikte)
$\{\emptyset,\{\emptyset\}\}$ olarak ele almamızdır.

Neden bir Sonsuzluk Aksiyomu'nu ötekine tercih edelim? Başlıca pratik
neden, von Neumann yaklaşımının sonlu ötesi sayıları ele alacak biçimde iyi
``ölçeklenmesidir''. Bunu \olref[ordinals][]{chap} bölümünden başlayarak
inceleyeceğiz. Bununla birlikte, yalnızca \emph{aritmetik yapma} açısından
iki yaklaşım da eşit ölçüde işe yarar. Öyleyse biri size doğal sayıların
küme \emph{olduğunu} söylerse açık soru şudur: \emph{Hangi kümeler?}

Bu kesin soruyu \citet{Benacerraf1965} ünlü hâle getirmiştir. Ancak bunun,
daha önce birçok kez karşılaştığımız bir olgunun yalnızca en tanınmış örneği
olduğunu vurgulamak gerekir. Temel nokta şudur: küme teorisi bize bir dizi
``sezgisel'' varlık türünü \emph{simüle etmenin} yolunu verir; evet,
gerçelleri, rasyonelleri, tam sayıları ve doğal sayıları, ama aynı zamanda
sıralı ikilileri, fonksiyonları ve bağıntıları da. Ne var ki küme teorisi
bize hiçbir zaman \emph{biricik} bir simülasyon seçeneği sunmaz. Bize aynı
ölçüde iyi hizmet edecek alternatifler \emph{her zaman} vardır.

\end{document}
